\documentclass[12pt,a4paper]{article}

\usepackage[utf8]{inputenc}
\usepackage[T1]{fontenc}
\usepackage{amsmath}
\usepackage{graphicx}
\usepackage{verbatim}
\usepackage{geometry}
\usepackage{hyperref}

\geometry{margin=1in}

\title{Anomaly Detection using Convolutional Autoencoders on Tiny-ImageNet}
\author{}
\date{}

\begin{document}

\maketitle

\section{Problem definition \& methodology}

The objective of this assignment was to investigate whether Convolutional Autoencoders (CAEs) can be used for anomaly detection on Tiny-ImageNet (10 classes). 

The task was to:
\begin{itemize}
    \item Train an autoencoder on 9 classes (seen data)
    \item Hold out 1 class as OOD (unseen)
    \item Compare reconstruction error (MSE) between seen and unseen samples
\end{itemize}

\textbf{Methodology:}
\begin{enumerate}
    \item Selected one class (e.g., class 3) as OOD
    \item Trained a convolutional autoencoder on remaining 9 classes
    \item Evaluated reconstruction error on:
    \begin{itemize}
        \item Seen classes
        \item Unseen (OOD) class
    \end{itemize}
    \item Visualized:
    \begin{itemize}
        \item Training loss
        \item Reconstruction error per class
        \item Original vs reconstructed images
    \end{itemize}
    \item Extended experiment by:
    \begin{itemize}
        \item Repeating training for each class as OOD (10 models)
        \item Computing mean and variance of reconstruction error per class
    \end{itemize}
\end{enumerate}



\section{Hyperparameter tuning \& architecture designs}

\textbf{Model:} Convolutional Autoencoder

\begin{itemize}
    \item Encoder: 3 Conv layers (stride=2), channels: 3 $\rightarrow$ 32 $\rightarrow$ 64 $\rightarrow$ 128
    \item Decoder: 3 ConvTranspose layers (mirroring encoder)
    \item Activation: ReLU (hidden), Tanh (output)
    \item No BatchNorm or Dropout
\end{itemize}

\textbf{Training setup:}
\begin{itemize}
    \item Loss: Mean Squared Error (MSE)
    \item Optimizer: Adam ($lr = 10^{-3}$)
    \item Batch size: 128
    \item Epochs: 30
\end{itemize}

\textbf{Observation:} Training loss converged smoothly, indicating stable learning.

\begin{figure}[h]
    \centering
    \includegraphics[width=0.6\linewidth]{loss.png}
    \caption{Training loss curve}
\end{figure}



\section{Results, visualizations, and interpretations}

\textbf{Single OOD experiment (Class 3):}

\begin{itemize}
    \item Seen class MSE: 0.0284
    \item Unseen class MSE: 0.0168
\end{itemize}

Interestingly, the unseen class did not always produce higher error, indicating that reconstruction alone may not reliably detect anomalies.

\textbf{Class-wise reconstruction error:}

\begin{verbatim}
Class 0: Avg MSE = 0.028132
Class 1: Avg MSE = 0.036198
Class 2: Avg MSE = 0.018450
Class 3: Avg MSE = 0.016844
Class 4: Avg MSE = 0.028276
Class 5: Avg MSE = 0.027768
Class 6: Avg MSE = 0.043647
Class 7: Avg MSE = 0.023515
Class 8: Avg MSE = 0.020019
Class 9: Avg MSE = 0.029609
\end{verbatim}

\begin{figure}[h]
    \centering
    \includegraphics[width=0.7\linewidth]{bar.png}
    \caption{Reconstruction error per class (OOD highlighted)}
\end{figure}

\textbf{Reconstruction visualization:}

\begin{figure}[h]
    \centering
    \includegraphics[width=\linewidth]{recon.png}
    \caption{Original vs reconstructed images}
\end{figure}

The reconstructions are generally good, even for the unseen class, suggesting the model learns shared visual features rather than class-specific patterns.



\section{The key findings and explanations}

\begin{itemize}
    \item Autoencoders do not reliably detect anomalies:  
    Reconstruction error for unseen class is not consistently higher.

    \item Error depends on feature similarity:  
    Classes sharing visual patterns (texture, color, structure) are reconstructed well even if unseen.

    \item Low variance across OOD runs:  
\end{itemize}

\begin{verbatim}
Class 0: Mean=0.043255, Var=0.000002
Class 1: Mean=0.054327, Var=0.000002
Class 2: Mean=0.030177, Var=0.000001
Class 3: Mean=0.028947, Var=0.000001
Class 4: Mean=0.039932, Var=0.000001
Class 5: Mean=0.038491, Var=0.000001
Class 6: Mean=0.063730, Var=0.000005
Class 7: Mean=0.035143, Var=0.000001
Class 8: Mean=0.029017, Var=0.000001
Class 9: Mean=0.048054, Var=0.000008
\end{verbatim}

\begin{figure}[h]
    \centering
    \includegraphics[width=0.7\linewidth]{variance.png}
    \caption{Mean $\pm$ Std reconstruction error across OOD runs}
\end{figure}

\textbf{Conclusion:}

\begin{itemize}
    \item Reconstruction error is not strongly dependent on which class is OOD
    \item CAEs learn a shared latent representation of visual features
    \item This limits their effectiveness for anomaly detection in complex datasets
\end{itemize}

\textbf{Final insight:}  
Anomaly detection using autoencoders works only when anomalies are sufficiently different in feature space, not merely unseen during training.

\end{document}
